\documentclass[letter,11pt]{article}

\usepackage[english]{babel}
\usepackage[T1]{fontenc}
%\usepackage[ansinew]{inputenc}
\usepackage{lmodern}	% font definition
\usepackage{datetime}
\usepackage{verbatim}

\usepackage{parskip}
\usepackage{graphicx}
\usepackage[top=2.2cm,bottom=2.1cm,right=1.9cm,left=1.9cm]{geometry}
\usepackage{pdflscape}

\usepackage[urw-garamond]{mathdesign}

\usepackage{tikz}

\usetikzlibrary{arrows,shapes,shapes.callouts,decorations.text}

\usepackage{url}
\usepackage[ps2pdf,breaklinks=true,bookmarks=true,bookmarksopen,bookmarksopenlevel=1,pdfpagelayout=OneColumn,pagebackref=true]{hyperref}
\usepackage{breakurl}

\hypersetup{
    bookmarks=true,         % show bookmarks bar?
    unicode=false,          % non-Latin characters in Acrobat’s bookmarks
    pdftoolbar=false,        % show Acrobat’s toolbar?
    pdfmenubar=true,        % show Acrobat’s menu?
    pdffitwindow=true,     % window fit to page when opened pdfstartview={FitV},    % fits the width of the page to the window
    pdftitle={Whole-Body Neural Recordings in a Month},    % title
    pdfauthor={David Dalrymple},     % author
    pdfcreator={David Dalrymple},   % creator of the document
    pdfproducer={David Dalrymple}, % producer of the document
    pdfnewwindow=true,      % links in new window
    colorlinks=true,       % false: boxed links; true: colored links
    linkcolor=blue,          % color of internal links
    citecolor=green,        % color of links to bibliography
    filecolor=magenta,      % color of file links
    urlcolor=cyan           % color of external links
}

\renewcommand*{\backref}[1]{}
\renewcommand*{\backrefalt}[4]{%
  \ifcase #1 %
(Not cited.)%
\or
(Cited on page #2.)%
\else
(Cited on pages #2.)%
\fi
}

\newcommand{\attrib}[1]{\nopagebreak{\raggedleft\footnotesize #1\par}}
\newcommand{\todo}[1]{\textcolor{lightgray}{\textit{<<#1>>}}}
\newcommand{\tbc}{\begin{center} \todo{to be completed} \end{center}}
\newcommand{\tbcsubsubsection}[1]{ \refstepcounter{subsubsection}%
  \subsubsection*{\thesubsubsection \quad #1} \tbc}
\newcommand{\cel}{{\em C. elegans}}

\setlength{\parskip}{0.3cm plus3mm minus1mm}
\setlength{\parindent}{0cm}

\begin{document}

\begin{center}
	\textsc{\LARGE {\bfseries Whole-Body Neural Recordings in a Month}}\\[2mm]
  \textit{An amibitious proposal by}  \textrm{David Dalrymple}\\[1mm]
	%{\small {\bfseries \color{red} Draft 3} compiled \currenttime, \today\\[0mm]}
	\rule[2mm]{0.66\textwidth}{0.25mm}\\[4mm]
\end{center}

\thispagestyle{empty}

As a reader of this document, you are already aware of my goal to simulate the
complete nervous system of \cel\ in software, based on new experimental data
enabled by a suite of technologies coming to be known as ``optical patch-clamping.''
By systematically perturbing neurons using light-driven channels and pumps,
then observing the response of the network using fluorescent voltage sensors,
it is my hope and belief that we can accurately infer circuit dynamics.

A initial exploratory step towards this goal is to exclude the perturbative
aspect. This would mean that no special algorithms would be needed to perform
the physical experiment, since all the data can be recorded and analyzed
post-hoc. Even this level of recording would be, as far as I know,
unprecedented, and very likely to be scientifically interesting. The voltage
sensor and associated imaging tools are already working in other model systems,
like {\em E. coli} and cardiomyocytes, so the major barrier to this milestone
is the pan-neuronal expression of the voltage sensor in \cel.

I'm going to be traveling for much of July and August. I know next to nothing
about transgenesis, but I'm pushing to reach this goal by July 1. I'd like your
feedback on whether this proposed timeline makes any sense:

\begin{center}
\begin{tikzpicture}[scale=0.45,yscale=-1,rounded corners=3pt,transform canvas={xshift=-1cm}]
  \draw [->] (0,0) -- (0,15);
  \foreach \y in {0,...,15}
    \draw (-4pt,\y) -- (4pt,\y);
    \node[rectangle callout, fill=yellow, opacity=.8, overlay, callout absolute pointer={(0,0)}] at (-8,0.5) {Today, {\bfseries June 7}};
    \node[rectangle callout, fill=yellow, opacity=.8, overlay, callout absolute pointer={(0,15)}] at (-8,14.5) {{\bfseries June 22}};
    \node[rectangle callout, fill=blue!50, opacity=.8, overlay, callout absolute pointer={(0,2.5)},callout pointer width=0.7cm] at (7,0.5) {Acquire \cel\ specimens};
    \node[rectangle callout, fill=blue!50, opacity=.8, overlay, callout absolute pointer={(0,3.3)}] at (10,2.5) {Acquire voltage sensor DNA};
    \node[rectangle callout, fill=red!60, opacity=.8, overlay, callout absolute pointer={(0,3.6)}] at (9,4.5) {Fuse voltage sensor to {\em unc-119} promoter};
    \node[rectangle callout, fill=red!60, opacity=.8, overlay, callout absolute pointer={(0,6.3)}] at (7,6.3) {PCR transgene};
    \node[rectangle callout, fill=red!60, opacity=.8, overlay, callout absolute pointer={(0,6.6)},callout pointer width=0.7cm] at (7,8.5) {Perform microinjections};
    \node[rectangle callout, fill=red!60, opacity=.8, overlay, callout absolute pointer={(0,10.6)}] at (7,11.0) {Screen};
    \node[rectangle callout, fill=blue!50, opacity=.8, overlay, callout absolute pointer={(0,14.6)}] at (7,14.0) {Initial recordings};
\end{tikzpicture}
\end{center}

\vspace{7cm}

It would also be very helpful if someone experienced with these procedures could
coach me through them; in addition, it's not clear in what laboratory they should
be performed.

Thank you for your time and assistance!

\end{document}
